\chapter{Differentiable Resampling and Optimal Transport}
\label{ch:bf-dpf-resampling-ot}

The previous chapter showed how a particle flow may be turned into a proposal-
corrected particle method.  The present chapter addresses a different
mathematical difficulty: the standard resampling step is discrete and therefore
not pathwise differentiable.  If one wants a differentiable particle method,
one must either abandon pathwise differentiation at the resampling step or
replace the resampling mechanism by a smooth relaxation.

This chapter studies the main relaxations relevant to the BayesFilter DPF
program.  The key point is that differentiability is not obtained for free.  It
is obtained by changing the resampling map, and therefore by changing the object
whose gradient is computed.  The chapter is organized to make that trade-off
explicit.  Its source spine is the differentiable-resampling literature
\citep{zhumurphyjonschkowski2020,corenflos2021differentiable} together with
the optimal-transport and Sinkhorn background in
\citet{villani2003topics}, \citet{cuturi2013sinkhorn}, and
\citet{peyre2019computational}.

\section{Objects, conventions, and source roles}
\label{sec:bf-dr-objects}

The chapter uses the following objects.  Naming them up front is important
because several objects below are sometimes all called ``resampling'' even
though they are different mathematical maps.

\begin{center}
\small
\begin{tabular}{p{0.25\textwidth} p{0.65\textwidth}}
\toprule
Object & Meaning \\
\midrule
$\hat\pi_t^N$ & weighted empirical measure before resampling \\
$\hat\pi_{t,\mathrm{eq}}^N$ & equal-weight empirical measure after a
resampling or transport step \\
$a_j$ & discrete ancestor index selected by a categorical resampling rule \\
$\alpha$ & soft-resampling interpolation parameter \\
$\Pi$ & transport coupling between source particles and equal target masses \\
$\mathcal U(w,u)$ & feasible set of nonnegative couplings with source marginal
$w$ and target marginal $u$ \\
$C$ & transport cost matrix \\
$H(\Pi)$ & entropy convention, defined below as
$-\sum_{ij}\Pi_{ij}(\log \Pi_{ij}-1)$ \\
$\Pi_\varepsilon^*$ & exact optimizer of the entropic OT problem for fixed
$\varepsilon>0$ \\
$a,b$ & Sinkhorn scaling vectors in the factorization
$\Pi=\diag(a)K_\varepsilon\diag(b)$ \\
$\tilde x_t^{(j)}$ & barycentric output particle after the coupling is converted
to a deterministic equal-weight cloud \\
$\Pi_{\varepsilon,K}$ & finite-iteration Sinkhorn approximation computed in
code \\
\bottomrule
\end{tabular}
\end{center}

The source roles are deliberately limited.  The differentiable-resampling
sources motivate the need for a smooth replacement of hard ancestor selection;
the OT sources justify the coupling formulation; the Sinkhorn sources justify
the entropic scaling form and numerical caveats.  None of these citations by
itself proves that a relaxed resampling operator preserves the original
categorical law, preserves a pseudo-marginal likelihood estimator, or supplies
an HMC-valid target.  Those statuses must be established by the local equations
and by later target analysis.

\section{The resampling bottleneck}
\label{sec:bf-dr-bottleneck}

At time $t$, suppose the particle system after propagation and weighting is
represented by the weighted empirical measure
\begin{equation}
\label{eq:bf-dr-weighted-empirical}
  \hat \pi_t^N(dx)
  = \sum_{i=1}^N w_t^{(i)}\,\delta_{x_t^{(i)}}(dx),
  \qquad
  w_t^{(i)} \ge 0,
  \quad
  \sum_{i=1}^N w_t^{(i)}=1.
\end{equation}
Classical resampling attempts to replace this weighted cloud by an equally
weighted empirical measure
\begin{equation}
\label{eq:bf-dr-equal-weight-empirical}
  \hat \pi_{t,\mathrm{eq}}^N(dx)
  = \frac{1}{N}\sum_{j=1}^N \delta_{\tilde x_t^{(j)}}(dx),
\end{equation}
while preserving the filtering information carried by the weighted cloud as well
as possible.

For standard multinomial resampling, the new support points are selected by
random ancestor indices
$a_j \sim \mathrm{Categorical}(w_t^{(1:N)})$ and then setting
$\tilde x_t^{(j)} = x_t^{(a_j)}$.  This map is perfectly natural as a sampling
operation, but it is not pathwise differentiable with respect to the weights.
That is the resampling bottleneck.

\section{Standard resampling as a discontinuous map}
\label{sec:bf-dr-discontinuous-map}

Let $U_j \sim \mathrm{Uniform}(0,1)$ and write the cumulative weights as
\[
  C_k = \sum_{i=1}^k w_t^{(i)}.
\]
Multinomial resampling chooses the ancestor index $a_j$ such that
\begin{equation}
\label{eq:bf-dr-categorical-selection}
  C_{a_j-1} < U_j \le C_{a_j}.
\end{equation}
Thus the map from weights to ancestors is piecewise constant.  Away from the
measure-zero boundaries where a uniform draw lands exactly on a cumulative
weight break, the selected index does not change under an infinitesimal weight
perturbation.  Consequently, the pathwise derivative of the resampled particle
with respect to the weights is zero almost everywhere.

The one-dimensional two-particle case makes the discontinuity concrete.  Let
$N=2$, $w=(p,1-p)$, and draw a fixed common random number $U\in(0,1)$.  The
first ancestor is selected when $U\le p$ and the second ancestor is selected
when $U>p$:
\begin{equation}
\label{eq:bf-dr-two-particle-discontinuity}
  a(U,p)
  =
  \begin{cases}
    1, & U\le p,\\
    2, & U>p.
  \end{cases}
\end{equation}
For any fixed $U$, the realized map $p\mapsto a(U,p)$ is a step function.  Its
classical derivative is zero away from $p=U$ and undefined at $p=U$.  Reusing
the same $U$ makes the discontinuity visible; redrawing $U$ changes the random
map rather than supplying a pathwise derivative of the realized map.

This is the precise sense in which standard resampling blocks a naive pathwise
autodiff chain.  The issue is not that gradients are numerically inconvenient.
It is that the ancestor-selection map itself is discontinuous.  Score-function
or common-random-number estimators can still be designed around the discrete
operation, but they are not the same as differentiating the realized ancestor
selection path.  BayesFilter's DPF question is specifically about the latter:
whether the filtering likelihood path supplied to an HMC or training routine is
a smooth value-gradient computation.

\section{Soft resampling}
\label{sec:bf-dr-soft-resampling}

A simple differentiable replacement is soft resampling.  Instead of selecting a
new support point only from the categorical ancestor, one interpolates between
the ancestor particle and the weighted mean of the cloud
\citep{zhumurphyjonschkowski2020}.  A representative form is
\begin{equation}
\label{eq:bf-dr-soft-resampling}
  \tilde x_t^{(j)}
  = \alpha x_t^{(a_j)} + (1-\alpha)\bar x_t,
  \qquad
  \bar x_t = \sum_{i=1}^N w_t^{(i)}x_t^{(i)},
\end{equation}
where $\alpha \in [0,1]$ is an interpolation parameter.

For the bias calculation below, take $x_i,\bar x\in\R^{n_x}$, weights
$w_i\ge0$ with $\sum_iw_i=1$, and a scalar test function
$\varphi:\R^{n_x}\to\R$ that is twice continuously differentiable on a
neighborhood containing the line segments between $x_i$ and $\bar x$.  The
gradient $\nabla\varphi(x_i)$ is treated as a column vector, so
$\nabla\varphi(x_i)'(x_i-\bar x)$ is a scalar directional derivative.

The key structural point is immediate:
\begin{itemize}
  \item if $\alpha=1$, one recovers standard multinomial resampling and loses
    pathwise differentiability;
  \item if $\alpha<1$, the interpolation term depends smoothly on the weighted
    mean, but the sampled ancestor $a_j$ is still a discrete random index unless
    the method also replaces or conditions on that draw.
\end{itemize}

The benefit is therefore a smoother surrogate path, not an exact pathwise
derivative of categorical resampling.  If the ancestor draw is held fixed, the
map is differentiable only on that conditional branch; if the ancestor draw is
also relaxed, the target object has changed further and must be named.  The
cost is bias.  The weighted mean is preserved because, conditional on the
current weighted cloud,
\begin{align}
\label{eq:bf-dr-soft-mean-preserved}
  \E[\tilde x_t^{(j)} \mid X,w]
  &=
  \alpha \sum_{i=1}^N w_i x_i
  +(1-\alpha)\bar x
  =
  \bar x .
\end{align}
Consequently every affine test function $\varphi(x)=c+d'x$ is preserved in
conditional expectation.  Nonlinear functionals of the cloud need not be
preserved, because the argument of the test function is changed before the
expectation is taken.  For a smooth scalar test function $\varphi$,
\begin{equation}
\label{eq:bf-dr-soft-test-expectation}
  \E[\varphi(\tilde x_t^{(j)}) \mid X,w]
  = \sum_i w_i
      \varphi\!\left(\alpha x_i+(1-\alpha)\bar x\right).
\end{equation}
A first-order Taylor expansion around $x_i$ gives the displayed bias.  To see
the sign and order explicitly, write
$d_i=(1-\alpha)(\bar x-x_i)$.  Then
\[
  \alpha x_i+(1-\alpha)\bar x=x_i+d_i
\]
and, for each $i$,
\[
  \varphi(x_i+d_i)
  =
  \varphi(x_i)
  + \nabla\varphi(x_i)'d_i
  + \frac{1}{2}d_i'\nabla^2\varphi(\xi_i)d_i
\]
for some point $\xi_i$ on the segment between $x_i$ and
$\alpha x_i+(1-\alpha)\bar x$.  If the Hessian is bounded on the segment
neighborhood, summing the remainders gives $O((1-\alpha)^2)$.  Therefore
\begin{equation}
\label{eq:bf-dr-soft-test-bias}
  \E[\varphi(\tilde x_t^{(j)}) \mid X,w]
  - \sum_i w_i\varphi(x_i)
  =
  -(1-\alpha)\sum_i w_i
  \nabla\varphi(x_i)'(x_i-\bar x)
  + O((1-\alpha)^2).
\end{equation}
Thus the linear test-function case is preserved, but nonlinear summaries are
generally shifted.  Hence soft resampling is naturally interpreted as a
differentiable surrogate rather than as an exact replacement for multinomial
resampling.  The shrinkage can also reduce cloud diversity: if
$\alpha<1$, each selected ancestor is pulled toward the same weighted mean.
This may stabilize gradients, but it changes the empirical distribution whose
downstream likelihood or loss is being differentiated.

\section{Resampling as a transport problem}
\label{sec:bf-dr-transport-view}

A more geometric way to formulate equal-weight resampling is to view it as a
transport problem.  The starting measure is the weighted empirical measure
$\hat \pi_t^N$ in \eqref{eq:bf-dr-weighted-empirical}; the target class is the
set of equal-weight empirical measures of size $N$
\citep{villani2003topics,reich2013nonparametric}.

This makes it natural to ask for a transport plan that moves the weighted cloud
to an equally weighted cloud while minimizing a transport cost.  In particular,
one may interpret equal-weighting as a projection problem in the Wasserstein
geometry of probability measures.  Standard resampling does not solve this
projection problem exactly; it samples ancestor locations from the original
support.  Transport-based resampling instead tries to produce an equally
weighted cloud through an explicit coupling between the weighted and uniform
measures.

With source atoms $x_i$, target weights $u_j=1/N$, and nonnegative cost
$C_{ij}$, the coupling feasible set is
\begin{equation}
\label{eq:bf-dr-coupling-set}
  \mathcal U(w,u)
  =
  \left\{
    \Pi\in\R_+^{N\times N}:
    \sum_{j=1}^N \Pi_{ij}=w_i,\ 
    \sum_{i=1}^N \Pi_{ij}=u_j
  \right\}.
\end{equation}
The unregularized coupling problem is
\begin{equation}
\label{eq:bf-dr-unregularized-ot}
  \Pi_0^*
  \in
  \arg\min_{\Pi \in \mathcal U(w,u)}
  \sum_{i=1}^N\sum_{j=1}^N \Pi_{ij} C_{ij}.
\end{equation}
This is an exact linear-programming object for the chosen cost and feasible set.
It is not a statement that categorical resampling has become differentiable, and
it is not a proof that the categorical resampling law has been preserved.
Categorical resampling samples support points from the source cloud according to
ancestor indices.  The OT formulation instead chooses a coupling with prescribed
marginals and then usually converts that coupling into a deterministic output
cloud.  The equality of total mass and target weights is exact within
\eqref{eq:bf-dr-coupling-set}; equivalence to categorical resampling is not
being claimed.

This transport view is mathematically useful because it makes the relaxation
explicit: a differentiable transport map is not the same object as the discrete
categorical selection map.  It also separates two levels that are often blurred
in implementation discussions: the mathematical transport problem one wishes to
solve, and the numerical or regularized surrogate one actually computes.

\section{Entropic optimal transport resampling}
\label{sec:bf-dr-entropic-ot}

Let $w=(w_1,\ldots,w_N)$ be the source weights and let
$u=(1/N,\ldots,1/N)$ be the target equal weights.  Let the transport cost matrix
be
\begin{equation}
\label{eq:bf-dr-cost-matrix}
  C_{ij} = \|x_t^{(i)} - z_j\|^2,
\end{equation}
where $z_j$ are target support points or, in the barycentric form, support
locations constructed from the transport plan itself.

Assume for the derivation that $w_i>0$, $u_j>0$, $\sum_iw_i=\sum_ju_j=1$, and
that all entries of $C\in\R^{N\times N}$ are finite.  A coupling
$\Pi\in\R_+^{N\times N}$ has row sums in $\R^N$ and column sums in $\R^N$.
The scaling vectors $a,b$ used below also lie in $\R_+^N$; consequently
$K_\varepsilon b$, $K_\varepsilon' a$, and the elementwise products in
\eqref{eq:bf-dr-sinkhorn-marginals} are conformable $N$-vectors.

Following the entropy-regularized OT construction used by
\citet{corenflos2021differentiable}, the entropic optimal transport problem is
\begin{equation}
\label{eq:bf-dr-eot-primal}
  \Pi_\varepsilon^*
  = \arg\min_{\Pi \in \mathcal U(w,u)}
      \langle \Pi, C \rangle - \varepsilon H(\Pi),
\end{equation}
where this chapter uses the convention
\begin{equation}
\label{eq:bf-dr-entropy-convention}
  H(\Pi)
  =
  -\sum_{i=1}^N\sum_{j=1}^N
    \Pi_{ij}(\log \Pi_{ij}-1),
  \qquad
  \Pi_{ij}>0.
\end{equation}
With this convention, minimizing
$\langle\Pi,C\rangle-\varepsilon H(\Pi)$ is equivalent, up to constants, to
minimizing
$\sum_{ij}\Pi_{ij}C_{ij}
 + \varepsilon\sum_{ij}\Pi_{ij}(\log\Pi_{ij}-1)$
over the coupling polytope.  For fixed positive marginals and
$\varepsilon>0$, the entropy term makes the objective strictly convex on the
positive part of the feasible set, giving a unique optimizer for the regularized
problem.  It also changes the mathematical object:
$\Pi_\varepsilon^*$ solves a regularized problem, whereas $\Pi_0^*$ solves
\eqref{eq:bf-dr-unregularized-ot}.

The scaling form follows from the Lagrange first-order conditions.  Introduce
row and column multipliers $\lambda_i$ and $\mu_j$ and form
\[
  \mathcal L(\Pi,\lambda,\mu)
  =
  \sum_{i,j}\Pi_{ij}C_{ij}
  + \varepsilon\sum_{i,j}\Pi_{ij}(\log\Pi_{ij}-1)
  + \sum_i\lambda_i\left(\sum_j\Pi_{ij}-w_i\right)
  + \sum_j\mu_j\left(\sum_i\Pi_{ij}-u_j\right).
\]
For an interior optimizer, differentiating with respect to $\Pi_{ij}$ gives
the stationarity condition
\begin{equation}
\label{eq:bf-dr-eot-stationarity}
  C_{ij}
  + \varepsilon \log \Pi_{ij}
  + \lambda_i + \mu_j
  = 0 .
\end{equation}
Therefore
\begin{equation}
\label{eq:bf-dr-gibbs-factorization}
  \Pi_{ij}
  =
  \exp(-\lambda_i/\varepsilon)
  \exp(-C_{ij}/\varepsilon)
  \exp(-\mu_j/\varepsilon).
\end{equation}
Absorbing the multiplier exponentials into positive scaling vectors gives the
Sinkhorn representation
\begin{equation}
\label{eq:bf-dr-sinkhorn-form}
  \Pi_\varepsilon^* = \diag(a) K_\varepsilon \diag(b),
  \qquad
  (K_\varepsilon)_{ij} = \exp(-C_{ij}/\varepsilon),
\end{equation}
with $a_i>0$ and $b_j>0$ chosen to satisfy the marginal constraints.
Substituting \eqref{eq:bf-dr-sinkhorn-form} into
\eqref{eq:bf-dr-coupling-set} gives the scaling equations
\begin{align}
\label{eq:bf-dr-sinkhorn-marginals}
  a \odot (K_\varepsilon b) &= w,\\
  b \odot (K_\varepsilon' a) &= u.
\end{align}
The classical Sinkhorn updates are therefore
\begin{align}
\label{eq:bf-dr-sinkhorn-updates}
  a^{(k+1)} &= w \oslash (K_\varepsilon b^{(k)}),\\
  b^{(k+1)} &= u \oslash (K_\varepsilon' a^{(k+1)}),
\end{align}
where $\odot$ and $\oslash$ denote elementwise multiplication and division.
These updates solve the exact EOT problem only in the limit of convergence and
only for the stated kernel, marginals, cost matrix, and regularization
\citep{cuturi2013sinkhorn}.

The resulting equally weighted cloud is obtained by barycentric projection.  In
general, a column with mass $u_j>0$ has conditional source weights
$\Pi_{ij}^*/u_j$, so its barycenter is
\[
  \tilde x_t^{(j)}
  =
  \frac{1}{u_j}\sum_{i=1}^N \Pi_{ij}^*x_t^{(i)}.
\]
For the equal-weight target $u_j=1/N$, this becomes
\begin{equation}
\label{eq:bf-dr-barycentric-map}
  \tilde x_t^{(j)} = N \sum_{i=1}^N \Pi_{ij}^* x_t^{(i)}.
\end{equation}
The factor $N$ appears because the target marginal is $u_j=1/N$, so
$N\Pi_{ij}^*$ are the conditional weights of source atoms contributing to the
$j$th output particle.  If $x_i\in\R^{n_x}$, then
$\tilde x_t^{(j)}\in\R^{n_x}$ and the full output cloud again has $N$ support
points with equal weights.  This barycentric step converts a coupling into a
deterministic support cloud.  It is a map derived from the coupling, not the
coupling itself, and not a categorical ancestor draw.

This map is differentiable with respect to the particles and weights only after
one specifies what is being differentiated: the exact EOT optimizer
$\Pi_\varepsilon^*$ via an implicit differentiation argument, or a finite
unrolled Sinkhorn computation.  These are related but not identical objects.
At this point the chapter must keep four layers distinct: the exact categorical
resampling law, the unregularized OT problem, the entropically regularized OT
problem, and the finite-iteration Sinkhorn solve used to approximate the
regularized OT optimum in practice.

\begin{center}
\scriptsize
\begin{tabular}{p{0.20\textwidth} p{0.25\textwidth} p{0.23\textwidth} p{0.22\textwidth}}
\toprule
Differentiated object & Value object & Derivative route & Non-claim \\
\midrule
Exact regularized EOT optimizer &
$\Pi_\varepsilon^\star$ and its barycentric projection for fixed
$C,w,u,\varepsilon$ &
Implicit differentiation of the optimality or scaling equations, under
regularity and positive-marginal conditions &
Not a derivative of categorical resampling and not an unregularized OT claim. \\
Finite unrolled Sinkhorn solver &
$\Pi_{\varepsilon,K}$ produced by exactly $K$ implemented iterations &
Autodiff through the implemented computational graph &
Not a derivative of the converged EOT optimizer unless convergence and
truncation error are separately controlled. \\
Implicit fixed-point solver &
A fixed point satisfying the stated Sinkhorn equations to a declared tolerance &
Linearized fixed-point or KKT system &
Not valid without the fixed-point regularity, residual, and linear-solve
conditions. \\
HMC or training scalar in code &
Whatever scalar value actually uses the transport layer &
Derivative of that same scalar only &
Not a gradient of the original filtering likelihood unless a correction or
equivalence proof is supplied. \\
\bottomrule
\end{tabular}
\end{center}

\section{Numerical-analysis obligations for Sinkhorn}
\label{sec:bf-dr-sinkhorn-numerics}

The limit $\varepsilon\downarrow0$ connects EOT heuristically to the
unregularized OT objective, but it is not a harmless numerical limit.  The
kernel entries
\[
  (K_\varepsilon)_{ij}=\exp(-C_{ij}/\varepsilon)
\]
can underflow when costs are large relative to $\varepsilon$.  Underflow changes
the effective support of the kernel and can make the scaling equations
\eqref{eq:bf-dr-sinkhorn-marginals} numerically singular.  Stabilized
implementations therefore often work in log-scaled dual variables rather than
forming all entries of $K_\varepsilon$ in floating point
\citep{schmitzer2019stabilized}.

In practice, code computes a finite-iteration plan
$\Pi_{\varepsilon,K}$ rather than the exact optimizer
$\Pi_\varepsilon^*$.  The relevant residuals are
\begin{align}
\label{eq:bf-dr-sinkhorn-residuals}
  r_{\mathrm{row}}
  &= \left\|\Pi_{\varepsilon,K}\mathbf 1 - w\right\|_1,\\
  r_{\mathrm{col}}
  &= \left\|\Pi_{\varepsilon,K}'\mathbf 1 - u\right\|_1 .
\end{align}
Here $\mathbf1\in\R^N$, so both residual arguments are $N$-vectors.  The
row-residual checks the source marginal and the column-residual checks the
equal-weight target marginal.
Small residuals are necessary for a credible transport computation, but they do
not remove entropic bias relative to \eqref{eq:bf-dr-unregularized-ot} or the
target change relative to categorical resampling.

The gradient path must also be named.  Unrolled differentiation treats the
$K$ Sinkhorn iterations as the computational graph.  It is transparent but
requires memory proportional to the number of stored iterations, unless
checkpointing or recomputation is used.  Implicit differentiation instead
differentiates the fixed-point or optimality equations; it can reduce memory
but requires solving a linearized system and is valid only under the regularity
conditions of that fixed-point representation.  Either way, the gradient is the
gradient of the selected numerical or regularized object, not automatically the
gradient of a categorical-resampling likelihood estimator.

The leading memory footprint of a dense Sinkhorn layer is $O(N^2)$ for the cost,
kernel, or plan matrices, and the naive runtime is $O(N^2K)$ for $K$ scaling
iterations.  Those costs matter in banking-scale state-space applications
because large particle counts, long time series, and repeated HMC evaluations
compound the expense.  A reviewer should therefore see the numerical route,
the residual tolerance, the stabilization method, and the gradient path before
accepting any claim about differentiable-resampling suitability.

\section{Bias versus differentiability}
\label{sec:bf-dr-bias-vs-diff}

The central mathematical trade-off of this chapter is now explicit.

\subsection{Standard multinomial resampling}

Standard multinomial resampling preserves the intended categorical resampling
law, but it is not pathwise differentiable because the selected ancestor index is
locally constant in the weights except at discontinuity boundaries.

\subsection{Soft resampling}

Soft resampling is pathwise differentiable through the mean-interpolation term,
but it changes the resampling map by shrinking particles toward the weighted
mean.  Equation~\eqref{eq:bf-dr-soft-test-bias} is the local expression of the
resulting nonlinear-functional bias.

\subsection{Entropic OT resampling}

Entropic OT resampling is differentiable through the Sinkhorn/transport solve,
but it replaces categorical resampling by a relaxed transport projection.  It
therefore introduces approximation through the entropic regularization and, in
practice, through finite Sinkhorn iteration budgets.

Thus differentiability is obtained by replacing the exact categorical resampling
map with a smoother object.  The smoother object may be a useful approximation,
but it is not the same map.  More precisely, one should distinguish:
\begin{enumerate}[label=(\roman*)]
  \item the exact categorical resampling law;
  \item the exact solution of the unregularized OT problem when such a transport
    interpretation is used;
  \item the entropically regularized OT problem actually chosen for smoothness;
  \item the finite-iteration Sinkhorn approximation used in code;
  \item the solver-differentiated objective defined by unrolling or implicitly
    differentiating the chosen Sinkhorn computation.
\end{enumerate}
This distinction is essential for later HMC interpretation: if one differentiates
through a relaxed resampling operator, one must be explicit about whether the
resulting chain is intended to target the original posterior, a relaxed
posterior, or a surrogate objective.

\begin{table}[ht]
\centering
\small
\begin{tabular}{p{0.17\textwidth} p{0.23\textwidth} p{0.22\textwidth} p{0.24\textwidth}}
\toprule
Layer & Mathematical object & Exact or approximate? & Source of approximation \\
\midrule
Categorical resampling & ancestor draw from $\mathrm{Categorical}(w)$ & exact for the classical resampling law & no pathwise derivative through realized ancestors \\
Unregularized OT & coupling $\Pi_0^*$ solving \eqref{eq:bf-dr-unregularized-ot} & exact for the chosen transport projection & it is already a transport formulation, not categorical resampling \\
Entropic OT & coupling $\Pi_\varepsilon^*$ solving \eqref{eq:bf-dr-eot-primal} & exact for the regularized problem & entropy changes the unregularized OT objective \\
Finite Sinkhorn & numerical plan $\Pi_{\varepsilon,K}$ & approximate & finite iterations, tolerance, log-domain stabilization, floating-point error \\
Solver-differentiated objective & scalar produced by differentiating a finite or implicit Sinkhorn path & exact only for the selected computational graph or fixed-point equations & gradient belongs to the relaxed numerical object, not automatically to categorical resampling \\
\bottomrule
\end{tabular}
\caption{OT object versus solver approximation.}
\label{tab:bf-dr-ot-object-solver}
\end{table}

\begin{table}[ht]
\centering
\small
\begin{tabular}{p{0.16\textwidth} p{0.19\textwidth} p{0.18\textwidth} p{0.20\textwidth} p{0.16\textwidth}}
\toprule
Method & Exactness status & Differentiability source & Bias source & Runtime issue \\
\midrule
Categorical & exact resampling law & none pathwise & none relative to the classical resampling law & cheap but discontinuous \\
Soft & surrogate law & smooth mean interpolation & shrinkage toward $\bar x$ for nonlinear functionals & cheap but diversity-reducing \\
EOT & exact for regularized OT & smooth Sinkhorn scaling & entropy relative to unregularized OT and categorical law & $O(N^2K)$ Sinkhorn work \\
Finite Sinkhorn & numerical approximation to EOT & unrolled or implicit differentiated iterations & solver residual and tolerance choice & stability, memory, iteration count \\
Solver gradient & selected computational objective & autodiff or implicit fixed-point derivative & mismatch if reported as original likelihood gradient & memory or linear-solve burden \\
\bottomrule
\end{tabular}
\caption{Resampling family comparison.}
\label{tab:bf-dr-family-comparison}
\end{table}

\section{Implementation debug map}
\label{sec:bf-dr-debug-map}

The main implementation failures are diagnostic of different mathematical
layers.  A gradient that vanishes at resampling points points to accidental hard
ancestor selection.  Marginal mismatch in the transport plan points to an
incomplete Sinkhorn solve.  Exploding or NaN gradients at small $\varepsilon$
usually point to unstabilized kernel scaling rather than to the OT formulation
itself.

\begin{table}[ht]
\centering
\small
\begin{tabular}{p{0.22\textwidth} p{0.32\textwidth} p{0.30\textwidth}}
\toprule
Problem & Likely mathematical cause & Relevant literature \\
\midrule
Zero pathwise gradient at resampling & hard categorical ancestor selection & differentiable resampling in \citet{zhumurphyjonschkowski2020} \\
Particle cloud collapses toward one point & soft-resampling shrinkage too strong & soft-resampling bias analysis; equation~\eqref{eq:bf-dr-soft-test-bias} \\
Sinkhorn rows or columns do not match desired marginals & finite-iteration residual or numerical underflow & \citet{cuturi2013sinkhorn,peyre2019computational} \\
NaN or unstable gradients for small $\varepsilon$ & Gibbs kernel underflow or unstable scaling & stabilized/log-domain Sinkhorn \citep{schmitzer2019stabilized} \\
Posterior shifts after adopting EOT & relaxed resampling target differs from categorical law & EOT resampling \citep{corenflos2021differentiable} and HMC target analysis in Chapter~\ref{ch:bf-dpf-hmc-target-suitability} \\
Gradient changes when iteration budget changes & differentiated scalar is tied to $\Pi_{\varepsilon,K}$ or to an implicit solver tolerance & numerical-analysis obligations in Section~\ref{sec:bf-dr-sinkhorn-numerics} \\
\bottomrule
\end{tabular}
\caption{Implementation debug map for differentiable resampling and OT.}
\label{tab:bf-dr-debug-map}
\end{table}

\begin{table}[ht]
\centering
\small
\begin{tabular}{p{0.18\textwidth} p{0.24\textwidth} p{0.24\textwidth} p{0.20\textwidth}}
\toprule
Item & Source claim & Local derivation task & Debugging relevance \\
\midrule
Soft resampling & differentiable relaxation of hard resampling & derive preserved mean and nonlinear bias in \eqref{eq:bf-dr-soft-test-bias} & explains shrinkage and diversity loss \\
OT projection & equal-weighting can be posed as a transport coupling & separate \eqref{eq:bf-dr-unregularized-ot} from categorical resampling & prevents exactness overclaims \\
EOT/Sinkhorn & entropy gives smooth scaling form & identify $\Pi_\varepsilon^*$ and barycentric map & locates regularization and solver errors \\
Finite solver & implementation computes $\Pi_{\varepsilon,K}$ & distinguish numerical residual from OT definition & gives tolerance and stability diagnostics \\
Solver differentiation & code differentiates an unrolled or implicit solver path & identify the differentiated scalar and residual contract & prevents gradient-object overclaims \\
\bottomrule
\end{tabular}
\caption{Source claim, local derivation, and debugging use for E3.}
\label{tab:bf-dr-source-claim-map}
\end{table}

\FloatBarrier

\section{Why learned or amortized OT comes later}
\label{sec:bf-dr-why-learned-later}

A learned or amortized OT operator is a further approximation layer on top of
the OT baseline developed above.  It does not replace the need to understand the
OT resampling map itself.  For this reason, the present chapter stops at the OT
and Sinkhorn level.  The next chapter will treat learned or amortized transport
operators explicitly as approximations to an already relaxed transport map.

\section{What this chapter does and does not claim}
\label{sec:bf-dr-chapter-boundary}

This chapter develops the mathematical structure of differentiable resampling
and transport-based equal-weighting.  It does \emph{not} yet claim:
\begin{itemize}
  \item that soft resampling preserves the exact categorical resampling law;
  \item that OT resampling is identical to standard multinomial resampling;
  \item that learned/amortized OT is already covered by the same mathematics as
    the exact OT relaxation;
  \item that finite Sinkhorn is a black-box differentiable layer whose gradient
    has the same status as the original categorical resampling likelihood;
  \item that any of these relaxed resampling schemes are automatically final
    HMC targets for nonlinear structural models.
\end{itemize}
Its role is narrower and prior to those questions: it makes explicit how
resampling is modified to obtain differentiability and where the resulting bias
or relaxation enters.
